% ID: SEL-D0075-OB-E01
% Tanda S075.
\textbf{a)} En el primer cuadrante tomamos los puntos situados por debajo de las tres rectas. En el eje vertical el límite activo es $y=3$ y en el horizontal es $x=5/2$. Las dos rectas decrecientes se cruzan en $(2,2)$, y $2x+y=6$ corta a $y=x+3$ en $(1,4)$. El cruce de $4x+y=10$ con $y=x+3$ es $(7/5,22/5)$, que incumple $2x+y\le6$. Los vértices en orden son $(0,0)$, $(5/2,0)$, $(2,2)$, $(1,4)$ y $(0,3)$.

\begin{center}
\setlength{\unitlength}{1mm}
\begin{picture}(105,63)(-8,-8)
\put(0,0){\vector(1,0){90}}\put(0,0){\vector(0,1){51}}
\put(92,0){$x$}\put(1,50){$y$}
{\color{gray}
\qbezier(0,48)(37.5,24)(75,0)
\qbezier(21.25,52.8)(41.875,26.4)(62.5,0)
\qbezier(0,24)(42.5,37.6)(85,51.2)}
\qbezier(0,0)(31.25,0)(62.5,0)
\qbezier(62.5,0)(56.25,8)(50,16)
\qbezier(50,16)(37.5,24)(25,32)
\qbezier(25,32)(12.5,28)(0,24)
\qbezier(0,24)(0,12)(0,0)
\put(-5,-5){\scriptsize $0$}
\put(57,-6){\scriptsize $(5/2,0)$}
\put(51,17){\scriptsize $(2,2)$}
\put(25,34){\scriptsize $(1,4)$}
\put(1,25){\scriptsize $(0,3)$}
\put(13,9){\scriptsize Región factible}
\end{picture}
\end{center}

\textbf{b)} Escribimos $f(x,y)=2(2x+y)-3$. Como $2x+y\le6$, tenemos $f(x,y)\le9$. La igualdad se da exactamente en la parte factible de $2x+y=6$, es decir, en todo el segmento cerrado de $(1,4)$ a $(2,2)$. El máximo es $\boxed{9}$.
